% The preamble of the knowledge base guide (D83), as a document must build with it on the builder (D84).
\documentclass[12pt,a4paper]{article}
\usepackage[margin=2cm,headheight=15pt]{geometry}
\usepackage{amsmath,amssymb,amsthm}
\usepackage{fontspec}
\setmainfont{CMU Serif}
\usepackage{polyglossia}
\setdefaultlanguage{russian}
\setotherlanguage{english}
\usepackage{fancyhdr}
\usepackage{enumitem}
\usepackage[hidelinks]{hyperref}

\setlength{\parindent}{0pt}
\setlength{\parskip}{0.4em}
\pagestyle{fancy}
\fancyhf{}
\fancyhead[L]{Курс}
\fancyhead[R]{\textit{Математический анализ, 2026/27}}
\fancyfoot[C]{\thepage}
\setlist[enumerate,1]{label=(\alph*)}

\providecommand{\tg}{\operatorname{tg}}
\providecommand{\arctg}{\operatorname{arctg}}
\newcommand{\task}[1]{\par\vspace{1.5em}\textbf{#1.}\ }
\newcommand{\sol}{\par\textbf{Решение.}\ }

\newtheorem{theorem}{Теорема}

\begin{document}
\begin{center}\textbf{\Large Семинар 1. Домашнее задание}\end{center}

\task{1} Найдите предел.
\begin{enumerate}
\item $\lim_{x\to0}\frac{\tg 3x}{x}$
\end{enumerate}
\sol
\begin{enumerate}
\item $\lim_{x\to0}\frac{\tg 3x}{x}=\lim_{x\to0}\frac{3x}{x}=3$.
\end{enumerate}

\begin{theorem}\label{t}Пример теоремы, \textit{курсив} и \textbf{\textit{жирный курсив}}.\end{theorem}
По теореме~\ref{t}.
\end{document}
